\documentclass[10pt]{article}

% graphicx package, useful for including eps, pdf, and raster graphics
\usepackage{graphicx}
\DeclareGraphicsExtensions{.png,.png,.jpg,.pdf}

% basic packages
\usepackage{fancyhdr}
\usepackage{color}
\usepackage{parskip}
\usepackage{float}
\usepackage{microtype}
\usepackage{url}
\usepackage{hyperref}
\usepackage{booktabs}
\usepackage{makecell}
\usepackage{multirow}
\usepackage{pbox}
\usepackage{enumitem}

% fonts
\usepackage{fontspec}
\setmainfont{DroidSans}[
  Extension   = .ttf,
  UprightFont = *,
  BoldFont    = *-Bold ]
\setmonofont{DroidSansMono}[
  Extension   = .ttf,
  UprightFont = * ]

\hypersetup{colorlinks,urlcolor=blue}
\urlstyle{same}

% page layout
\usepackage[top=1in, bottom=1in, left=0.75in, right=0.75in, includefoot, heightrounded]{geometry}

% Useful for seeing the layout as you tweak it.
%\usepackage{showframe}

\setlength{\headheight}{1.1in}
\pagestyle{fancy}
\fancyhf{}

\lhead{\pbox[b][1in][c]{\textwidth}{
    \small
    \textbf{Prueba realizada por}\\
    Northwest Genomics Center\\
    University of Washington School of Medicine\\
    Genome Sciences Box 355065\\
    3720 15th Ave NE\\
    Seattle, WA 98195-5065\\
    206-616-5859}}
\rhead{\includegraphics[height=0.75in,trim=0 -25 0 0]{logo}}

\newcommand{\PageLine}{\rule{\textwidth}{0.25mm}}

\begin{document}
\sloppy % prevent overfilled lines which occur with long URLs

\begin{description}[font=\normalfont,align=left,labelwidth=18em]
\item [\textbf{\VAR{pat_name|e}}]
\item [\textbf{Fecha de nacimiento}] \VAR{birth_date|e}
\item [\textbf{Muestra}]
  %- if swab_type == "ans"
  Hisopo nasal anterior,
  %- elif swab_type == "mtb"
  Hisopo nasal del cornete medio,
  %- else
  \errmessage{Invalid swab_type "\VAR{swab_type}"}
  %- endif

  extraída por el paciente

\item [\textbf{Código de barras de la extracción}] \VAR{qrcode|e}
\item [\textbf{Fecha de la extracción}] \VAR{collect_ts|e}
\item [\textbf{Fecha del informe}]
  %- if status_code == "never-tested"
  Nunca probado
  %- else
  \VAR{result_ts|e}
  %- endif
\end{description}

\PageLine

\begin{center}
\Large
\textbf{RESULTADO DE LA INVESTIGACIÓN: le recomendamos que informe este
  resultado al proveedor de atención médica.}
\end{center}

\bigskip

\large \underline{\textbf{Acerca de esta prueba}}

\begin{itemize}
\item

  COVID-19 es una \textbf{enfermedad} provocada por un \textbf{virus} denominado
  SARS-CoV-2. Esta prueba de investigación detecta el virus en la muestra que se
  extrajo de la nariz para determinar si tiene COVID-19. El término técnico para
  este tipo de prueba es ``reacción en cadena de la polimerasa con transcripción
  inversa'', o ``RT-PCR'' para abreviar.

\end{itemize}

\bigskip

\large \underline{\textbf{Recordatorios importantes}}

Antes de consultar el resultado de la prueba de investigación, debe considerar
lo siguiente:

\begin{itemize}
\item

  El \textbf{tipo de muestra} que extrajo y \textbf{las pruebas en el
  laboratorio} forman parte de un \textbf{estudio de investigación} y este
  informe describe el resultado de la \textbf{prueba de investigación}.

\item
  %- if staff_observed == "f"
  Extrajo esta muestra y la etiquetó usted mismo (o para otra persona) y la
  devolvió al laboratorio mediante un servicio de entrega. No la extrajo,
  etiquetó o transportó un profesional de la salud con licencia, como un
  enfermero o médico.

  %- elif staff_observed == "t"
  Usted recogió esta muestra por su cuenta (o para otra persona) y fue etiquetada
  y devuelta al laboratorio a través de un servicio de mensajería o de otro
  tipo de profesional que no se dedica a la atención médica.  No la extrajo o
  transportó un profesional de la salud con licencia, como un enfermero o
  médico.

  %- else
  \errmessage{Unknown value: staff_observed = "\VAR{staff_observed}"}

  %- endif

  Esto podrá afectar el resultado de la prueba, y es posible que un proveedor de
  atención médica desee volver a realizarla.

\item

  Participar en el estudio de investigación no reemplaza la atención médica
  habitual. Si se siente muy mal o está preocupado por el resultado de la
  prueba, se debe comunicar con el médico o proveedor de atención médica.

\end{itemize}

\bigskip
\PageLine

\large \underline{\textbf{Resultado de la prueba de investigación del SARS-CoV-2
  (COVID-19)}}

%- if status_code == "negative"
No se detectó (negativo)\\
%- elif status_code == "positive"
Se detectó (positivo)\\
%- elif status_code == "inconclusive"
No concluyente\\
%- elif status_code == "never-tested"
No se pudo analizar\\
%- else
\VAR{0/0} %# poor man's raise
%- endif

\PageLine
\bigskip

\large \underline{\textbf{Entender el resultado}}

Esta tabla proporciona todos los resultados que la prueba puede revelar y
describe lo que significa cada uno de ellos. Puede imprimir o enviar por correo
electrónico el resultado de esta \textbf{prueba de investigación} para
informárselo al médico o proveedor de atención médica.

Sin importar el resultado de la prueba, debe quedarse en su casa y evitar estar
cerca de otras personas si se siente mal. \textbf{Si sus síntomas continúan, se
siente muy mal, está preocupado por la enfermedad o tiene preguntas sobre el
resultado de la prueba de investigación, debe comunicarse con el médico o
proveedor de atención médica.}

\begin{itemize}


\item

  \textbf{Si la prueba indica ``No se detectó (negativo)''}, significa que nuestro
  laboratorio no encontró el virus que provoca la COVID-19 en la muestra que nos
  envió. Las pruebas no son perfectas, y existe una pequeña posibilidad de que
  tenga el virus a pesar de que el resultado de la prueba haya sido negativo.


\item

  \textbf{Si la prueba indica ``Se detectó (positivo)''}, significa que nuestro
  laboratorio encontró el virus que provoca la COVID-19 en la muestra que nos
  envió. Debe quedarse en su casa y evitar estar cerca de otras personas,
  incluso si no tiene síntomas de COVID-19. \textbf{No vaya a trabajar ni asista
  a la escuela.}

  Si no ha recibido una llamada de un representante de nuestro programa, llame
  al 206-616-5859 para analizar su resultado.

  El resultado se compartirá con el Estado de Washington y los departamentos de
  salud locales. \textbf{No} es necesario que usted informe estos resultados. El
  departamento de salud puede llamarlo para hacerle preguntas sobre la
  enfermedad y las personas con las que pudo haber estado mientras estuvo
  enfermo.


\item

  \textbf{Si la prueba indica ``No concluyente''}, significa que nuestro
  laboratorio no está seguro del resultado de la prueba. Las pruebas no son
  perfectas, y el resultado no concluyente ocurre cuando se encuentra una
  pequeña cantidad de SARS-CoV-2 en la muestra, pero no lo suficiente como para
  que podamos afirmar que tiene COVID-19.

  Las pautas actuales de salud pública sugieren que se haga una prueba
  adicional, pero debe asumir que este resultado es positivo y quedarse en casa,
  alejado de los demás, incluso si no tiene síntomas de COVID-19. \textbf{No
  vaya a trabajar ni asista a la escuela.}

  Si no ha recibido una llamada de un representante de nuestro programa, llame
  al 206-616-5859 para analizar su resultado.

\item

  \textbf{Si la prueba indica ``No se pudo analizar''}, significa que no había
  suficiente contenido en la muestra como para que la pudiéramos examinar, o se
  produjo un error cuando realizamos la prueba. Como no logramos analizar la
  muestra, no pudimos determinar si tiene o no el virus que provoca la COVID-19.


\end{itemize}

\bigskip

\large \underline{\textbf{Recursos}}

Si desea obtener más información sobre la COVID-19, consulte estos recursos de
los organismos de salud pública:

\begin{itemize}
\item
  Public Health \textemdash Seattle \& King County (Salud Pública \textemdash Seattle y el Condado de King)\\
  \url{https://www.kingcounty.gov/depts/health/covid-19/languages/spanish.aspx}

\item

  Washington State Department of Health (Departamento de Salud del Estado de Washington)\\
  \url{https://www.doh.wa.gov/Emergencies/NovelCoronavirusOutbreak2020COVID19/CoronavirusSpanish}

\item
  U.S. Centers for Disease Control and Prevention (CDC, Centros para el Control
  y la Prevención de Enfermedades de los Estados Unidos)\\
  \url{https://espanol.cdc.gov/coronavirus/2019-ncov/index.html}

\item
  Página web de recursos relacionados a la COVID-19 de la University of
  Washington\\
  \url{https://www.washington.edu/coronavirus/}

\end{itemize}

Si desea hablar por teléfono con alguien, comuníquese de alguna de las
siguientes maneras: 

\begin{itemize}

\item

  \textbf{Para realizar consultas médicas sobre la COVID-19:} comuníquese con el
  King County Call Center (Centro de Atención del Condado de King), entre las
  8:00 a. m. y las 7:00 p. m., llamando al \textbf{206-477-3977}

\item

  \textbf{Para realizar consultas no médicas sobre la COVID-19, lo que incluye
  asuntos relacionados con el cumplimiento y los negocios:} comuníquese con la
  King County COVID-19 Business and Community Information Line (línea de
  información comunitaria y comercial sobre la COVID-19 del Condado de King), de
  lunes a viernes entre las 8:30 a. m. y las 4:30 p. m., llamando al
  \textbf{206-296-1608}

\item

  \textbf{Preguntas generales sobre nuestro programa, incluidas preguntas sobre
  sus resultados:} Comuníquese con nuestro equipo llamando al
  \textbf{206-616-5859} o enviando un correo electrónico a
  \textbf{participanthelp@uw.edu}.

\item

  \textbf{Para realizar consultas generales sobre la COVID-19 en el Estado de
  Washington:} comuníquese con el Washington State Novel Coronavirus Call Center
  (Centro de Atención del Estado de Washington para el Nuevo Coronavirus), entre
  las 6:00 a. m. y las 10:00 p. m., llamando al \textbf{1-800-525-0127}

\item

  \textbf{Preguntas generales sobre la COVID-19 de la University of Washington:}
  Comuníquese con el equipo de respuesta a la COVID-19 del Environmental Health
  and Safety Department (EH\&S, Departamento de Salud y Seguridad Ambiental) de
  la University of Washington (UW) llamando al \textbf{206-616-3344} o enviando
  un correo electrónico a \textbf{covidehc@uw.edu}.

\end{itemize}

\bigskip

\large \underline{\textbf{Métodos y limitaciones}}

\textbf{Métodos}\\
Este ensayo se realiza para detectar SARs-CoV-2 utilizando la técnica PCR con
transcripción inversa. Este ensayo utiliza la amplificación de los segmentos de
virus ORf1B y S, y también detecta el control de amplificación de la RnaseP
humana. El University of Washington’s Northwest Genomics Center desarrolló esta
prueba y determinó sus características de rendimiento. Este laboratorio está
autorizado por el WA Department of Health como centro de pruebas médicas para
realizar pruebas de laboratorio clínico de alta complejidad.

\textbf{Limitaciones}

\begin{itemize}

\item

  Se brindan estos resultados como parte de un estudio de investigación.

\item

  %- if staff_observed == "f"
  Extrajo esta muestra y la etiquetó usted mismo (o para otra persona) y la
  devolvió al laboratorio mediante un servicio de entrega. No la extrajo,
  etiquetó o transportó un profesional de la salud, como un enfermero o médico.

  %- elif staff_observed == "t"
  Usted recogió esta muestra por su cuenta (o para otra persona) y fue etiquetada
  y devuelta al laboratorio a través de un servicio de mensajería o de otro
  tipo de profesional que no se dedica a la atención médica.  No la extrajo o
  transportó un profesional de la salud con licencia, como un enfermero o
  médico.

  %- else
  \errmessage{Unknown value: staff_observed = "\VAR{staff_observed}"}

  %- endif

  Esto podrá afectar el resultado de la prueba, y es posible que un proveedor de
  atención médica desee volver a realizarla.

\item

  Como este informe se basa en datos de un estudio de investigación, es posible
  que los resultados no puedan diagnosticar COVID-19, la enfermedad provocada
  por la infección con el virus SARS-CoV-2.

\end{itemize}

\end{document}
